\section{The original late-family setting and the local Hecke estimates}
\label{sec:global-local-estimates}

\status{This chapter reconstructs the numerical setting and the finite local
arguments on the original objects.  A supplied family, its original forms,
and their same-object ties are retained throughout.  The source is the
pinned commit recorded in the accompanying provenance ledger.}

\paragraph{Objects and conventions.}
Fix a source member $D$, a word $u$, and write
$K=D.\mathrm{what}_2.\mathrm{LL}$, $k=D.R.k=N(0)$,
$m=D.R.m$, $N=N(m)$, $f=D.R.f$, and $p=(D.\mathrm{mem}\,u).A.T$.
The prime $p$ is the original word prime.  At a coefficient place $w\mid p$,
write $\mathcal O_w$ for the original DVR and $e_w=v_w(p)$.
An original level-zero forms record is denoted $\mathcal F$; this avoids
using $T$ for both that record and the prime.  In the source its type is
\texttt{Forms0 D u}.  Its maps \texttt{atT}, \texttt{cls0}, \texttt{val0},
\texttt{lv0}, \texttt{up\_natural\_level\_zero},
\texttt{down\_natural\_level\_zero}, and \texttt{bw0} are the original
ones, obtained by projecting the supplied symmetric forms record.

\subsection{A direct size proof for the actual weights}

% BEGIN GLE1
\paragraph{GLE1: the original parameter windows give a common weight interval.}
\label{claim:GLE1}
Let $L:\texttt{SizeData}_2(D)$ and let $P:\texttt{ParameterChoice}'(D,L)$
be the supplied parameter record.  Put
\[
 \begin{gathered}
 x=L.\mathrm{sc}.x,\quad b=e^{x^\varepsilon},\quad a=xb,
 \quad X'=e^{x^{\varepsilon/3}},\quad b'=a/X',\\
 w''=b'/(xn),\quad S=L.\mathrm{Sc},\quad \eta=L.\mathrm{etaSz}.
 \end{gathered}
 \tag{GLE.1}
\]
Assume $x\ge4$ and $X'\ge x$.  Then, for every original word and every
embedding, every entry of every original lifted weight with $0\le j\le m$
satisfies
\[
 -2b\le\lambda_{j,\tau}(r)\le2xb.                 \tag{GLE.2}
\]
Moreover,
\[
 b\ge24x,\qquad b\ge3N,\qquad N\le b/4,
 \qquad a_{\max}\le2xb,
 \qquad b_{\mathrm{top},\tau}>0\quad(\tau\in\mathrm{tops}(D)).
 \tag{GLE.3}
\]

Here is the complete extraction from the parameter fields.  Since
$h_0=\eta/m$, the ninth inequality of \texttt{h0\_ge} gives
$Nm^{11}\le\eta$.  The field \texttt{mteta} gives
$m(t+1)\eta\le w''$, and \texttt{kS\_le} gives $kS\le w''$.
The rank laws give $m\ge2$, $k\ge5$, $n\ge1$, $N\ge1$, $t\ge1$;
\texttt{t\_le\_eta} gives $\eta\ge1$.  Consequently
\[
 Nm^{12}\le w'',\qquad 1\le\eta\le w'',\qquad
 0<S\le w'',\qquad b'=xnw'',\qquad 0<b'\le b.
 \tag{GLE.4}
\]
The last inequality follows directly from $b'=xb/X'$.
In particular $w''\ge4096N$, $b'\ge4096xN$, and
$w''\le b'/x\le b'/4$.  These imply the first three bounds in (GLE.3).

We check the first $k$ entries before the appended blocks.  Write
$u_r$ for the original half-integral Harish--Chandra entries and
\[
 \lambda_r=u_r-\bigl((k-1)/2-r\bigr).
 \tag{GLE.5}
\]
The intervening entries are integral and strictly interlace the $u_r$.
Thus $u_r-u_{r+1}\ge1$, and (GLE.5) is decreasing.  At a top place the
first and last entries are exactly $a_T$ and $-b_T$.  The original
window for $a_T$ gives $a\le a_T\le2a$.
For $b_T$, put $q=N_*x^\varepsilon L.\mathrm{LaG}/k$.
The fields \texttt{S\_le}, $S>0$, and $b'>0$ imply
$L.\mathrm{LaG}>0$.  The definition of \texttt{rhoP} is
\[
 \rho_P=4q+C_\rho\log x\,L.\mathrm{LaG}/k.
\]
All terms on the right are nonnegative, and
$\rho_P\le x^{-2}\le1$.  Hence $q\le1/4$ and the actual
\texttt{bT\_le} and \texttt{bT\_ge} fields give
\[
 b'\le b_T\le\tfrac54 b'.                       \tag{GLE.6}
\]
Every middle top entry belongs to $[0,S]$.  Indeed $u_1\le S$ and
the shift $(k-1)/2-1$ is nonnegative, so $\lambda_1\le S$.
At the other endpoint $u_{k-2}\ge0$ and the shift
$(k-1)/2-(k-2)$ is nonpositive, so $\lambda_{k-2}\ge0$.
The monotonicity of (GLE.5) places all intermediate entries between
these two bounds.  Off the top set, \texttt{uk\_extreme} and the
last cloud bound give $\lambda_{k-1}\ge0$; the first cloud bound and
\texttt{cloud\_place\_bound} give $\lambda_0\le3S$.
Thus all off-top entries belong to $[0,3S]$.
Since $3S\le3b'/4\le b\le a$, these entries also satisfy (GLE.2).
The maximum defining $a_{\max}$ is the maximum of $0$ and the original
top entries $a_T$; hence $a_{\max}\le2a$.  The original definition
$b_{\mathrm{top},\tau}=-\lambda_{0,\tau}(k-1)$ identifies it with
$b_T$ at every top place, proving its positivity.

For the appended block indexed by $0\le r<m$, the original weight
formula is $-w_r$, where
\[
 w_r=w_{r+1}^{\mathrm{raw}}-\delta_r,
 \qquad
 \delta_r=\frac{k}{2}+\frac{(2r+1)(t+1)}2,
 \qquad 0\le\delta_r\le N.                       \tag{GLE.7}
\]
The three raw windows are
\[
 \begin{array}{c|c}
 \text{top}&[b_T+b'/x,\ b_T+b'/x+\eta]\\
 \text{thick}&[b,\ b+\eta+1]\\
 \text{thin}&[w'',\ w''+\eta].
 \end{array}                                      \tag{GLE.8}
\]
Every lower endpoint is at least $w''\ge N$.
The top upper endpoint is at most $7b'/4$, the thick one at most
$b+b'/2$, and the thin one at most $b'/2$; here $1\le w''\le b'/4$.
Subtracting (GLE.7) therefore gives $0\le w_r\le2b$.
The source's literal \texttt{lam\_eq} identifies the first $k$ entries
with (GLE.5) and every remaining entry with one of these $-w_r$.
This proves (GLE.2) for all $j$, including $j=0$ and $j=m$.
No general \texttt{NumsL} conclusion is being presumed in this proof.
% END GLE1

% BEGIN GLE2
\paragraph{GLE2: \texttt{TuSize} supplies every clause of the eventual setting.}
\label{claim:GLE2}
For every original family $\mathrm{Fm}:\texttt{FamP}\,\varepsilon$,
\[
 \texttt{TuSize}(\mathrm{Fm})
 \quad\Longrightarrow\quad
 \texttt{SettingEventually}(\mathrm{Fm}).          \tag{GLE.9}
\]
The conclusion means that there is one $A$ such that, for every
$x\ge A$ with $x>\mathrm{Fm}.x_1$, every $i\in\mathrm{Fm}.I_x$ and
every word $u$ of $D=\mathrm{Fm}.D(x,i)$ satisfy all four original
\texttt{HidaSetting} requirements.  The threshold is chosen before
$i$ and $u$.

To check its existence directly, put $\delta=\varepsilon/3>0$.
For $x\ge1$, $\log x\le x^{\delta/2}/(\delta/2)$, by
$\log y\le y$ with $y=x^{\delta/2}$.
Thus $x^{\delta/2}\ge2/\delta$ implies $x\le e^{x^\delta}=X'$.
Increase a single family threshold until this holds, $x\ge32$, and
the \texttt{TuSize} threshold has been reached.  All original parameter
and size records are then supplied by \texttt{Fm.params} and
\texttt{Fm.sizes}; their scale satisfies $L.\mathrm{sc}.x=x$ and
$L.\mathrm{sc}.\varepsilon=\varepsilon$.

The lower $m$-size field says
\[
 m\ge x^{-\varepsilon/100}\sqrt{x\,x^{-\varepsilon/8}}
   =x^{1/2-29\varepsilon/400}\ge x^{2/5}\ge4.
 \tag{GLE.10}
\]
The middle inequality uses the original $0<\varepsilon<1/52$.
The lower prime bound in \texttt{TuSize}, namely
$p\ge b^{m/2}$, now gives $p\ge b^2$.
The original rank law is $f\le N$, so (GLE.3) yields
\[
 2fN+1\le2N^2+1<9N^2\le b^2\le p,
 \qquad 4a_{\max}\le8xb<24xb\le b^2\le p.
 \tag{GLE.11}
\]
For every $j\le m$ one has $N(j)\le N$.  By (GLE.2),
\[
 \begin{split}
 2\bigl(\lambda_{j,\tau}(r)-\lambda_{j,\tau}(s)+N(j)\bigr)
 &\le4xb+4b+2N\\
 &\le6xb<24xb\le b^2\le p.
 \end{split}                                      \tag{GLE.12}
\]
Indeed $N\le b/4$ and $x\ge4$ give $4b+2N\le2xb$.
This is the literal strict \texttt{Alcove} inequality, for every ordered
pair of entries.  The last clause $b_{\mathrm{top},\tau}\ge0$ follows
from GLE1.  Equations (GLE.11)--(GLE.12) prove the original setting,
including its strict inequalities, on one common threshold.

This proof needs only the lower \texttt{TuGe} part of \texttt{TuSize}.
The upper bound for $\log p$ and the logarithmic derivative estimate
are not extra premises.  They occur in the source's stronger general
$\texttt{NumsL}/\texttt{WtsL}$ supplier, whose remaining conclusions
are not claimed here.
% END GLE2

\subsection{The exact polynomial coordinate lattices}

% BEGIN GLE3
\paragraph{GLE3: integral and rational changes of polynomial coordinates.}
\label{claim:GLE3}
For a dominant integral weight $\lambda$ of rank $k$, put
$s=\lambda_{k-1}$ and $\nu_r=\lambda_r-s$.  Let
$\Delta_{\nu,\mathbb Z}$ be the original \texttt{delSq} lattice: the
integer span of all divided-power root words applied to the product of
principal minors of weight $\nu$.  For a coefficient DVR $\mathcal O_w$
put
\[
 \Lambda_w=\bigotimes_\tau
       (\mathcal O_w\otimes_{\mathbb Z}\Delta_{\nu_\tau,\mathbb Z}),
 \qquad V_K=K\Lambda_w.
 \tag{GLE.13}
\]
The tensor variables for different $\tau$ are disjoint.  After the
stored embedding $K\hookrightarrow\mathbb C$, these are exactly
\texttt{IntSpan D w (lam 0)} and \texttt{RatSpan D (lam 0)}.
Indeed the definitions are the additive spans of precisely the pure
tensors in (GLE.13), with respectively integral-at-$w$ and arbitrary
$K$ coefficients.  Negative determinant shifts belong to the value
function, not to this polynomial lattice.

For an invertible matrix $C$, use the source convention
\[
 R_\lambda(C)P(Z)=\det(C)^s P(CZ).                 \tag{GLE.14}
\]
If $C\in\mathrm{GL}_k(\mathcal O_w)$, then (GLE.14) preserves the
original integral lattice, with equality; if $C\in\mathrm{GL}_k(K)$,
it preserves its rational span.  To prove this without changing
lattices, a transvection $I+cE_{ab}$ acts by the finite Taylor expansion
$\sum_{r\ge0}c^r E_{ab}^{[r]}$.  Every divided-power root operator
preserves the span of root words.  A diagonal unit preserves each
integer row-weight summand, and its determinant shift is a unit even
when $s<0$.  Gaussian elimination over a local ring factors every
invertible matrix into transvections, permutations and diagonal units:
a unit determinant forces a unit entry in the first pivot column after
a row permutation; eliminate that column and induct.  Row permutations
are products of transvections and diagonal units.  This proves the
assertion and its inverse.  The same argument over $K$ proves the
rational assertion.  Tensoring proves both assertions for independent
matrices at the different $\tau$.

The lattice $\Delta_{\nu,\mathbb Z}$ has a basis homogeneous in all row
degrees.  Every generating root word has such a degree, it is a
submodule of the finite free integer module of polynomials of the
fixed total degree, and its distinct row-degree pieces have zero
intersection.  Choose an integer basis of each piece and take their
union.  If $r_\lambda$ is its rank and $d_j(h)$ is the $j$th row degree
of the basis element $h$, row permutations give equal total degrees
in all $k$ rows.  The total degree of every $h$ is $|\nu|$.  Hence
\[
 \sum_h(d_j(h)+s)=\frac{r_\lambda}{k}|\lambda|,
 \qquad d_j(h)+s\ge s.                            \tag{GLE.15}
\]
This elementary degree identity, including the shift $s$, will compute
the local length below.  It needs no semisimplicity in the residue
characteristic.

Finally, the polynomial-to-value map in the source is injective on
these coordinates.  Its frames range over all invertible matrices,
and the determinant factors in \texttt{fnOf} are nonzero.  A polynomial
vanishing on all such tuples is zero: for any fixed tuple of matrices
replace each component $Z$ by $Z+tI$.  Its determinant is nonzero for
all but finitely many $t$.  Here a separate variable $t_\tau$ is used
for each component.  The resulting polynomial vanishes on a product
of cofinite sets, hence is zero by induction on the number of
variables.  Setting all $t=0$ proves the assertion on every tuple.
The ordinary polynomial identity principle over the infinite field
$\mathbb C$ finishes the proof.  Therefore the displayed coordinate
transports are equalities of the actual forms, not merely maps on a
larger ambient polynomial space.
% END GLE3

% BEGIN GLE4
\paragraph{GLE4: exact coordinate descriptions of the original local lattices.}
\label{claim:GLE4}
Assume the original level-zero stabilizers are trivial, as in
\texttt{admS}; RSL7 supplies this from the actual forms in the selected
application.  Let $\mathcal C$ be the original finite class set.
For each $c\in\mathcal C$ choose a $p$-local basis $B_c$ of the original
dual class lattice
$H_c=\texttt{Ld D u 0 c}$.  The prescribed frame coordinates give a
$K$-linear isomorphism
\[
 V= ((D.\mathrm{fm}\,u).\mathrm{lat}\,0).V
       \ \simeq\ \bigoplus_{c\in\mathcal C}V_K.
 \tag{GLE.16}
\]
There are no further classwise constraints: the original
\texttt{cls0.stabI\_range} identifies the compatibility group at $c$
with its trivial stabilizer.  Arbitrary values at the preferred
representatives therefore extend uniquely by the original adelic
equivariance rule.  The \texttt{val0.rat} equality identifies those
whose coordinates lie in the original rational polynomial span.
Injectivity is GLE3 and the original injective embedding of $V$.
Thus both directions of (GLE.16) use the original rational form.

Write $S_w=(D.\mathrm{fm}\,u).S(0,w)$.  In (GLE.16) one has exactly
\[
                         S_w=\bigoplus_{c\in\mathcal C}\Lambda_w.
 \tag{GLE.17}
\]
First, two allowed local bases differ, after each relevant embedding,
by a matrix in $\mathrm{GL}_k(\mathcal O_w)$.  GLE3 therefore converts
the source's condition ``in every local basis'' into membership in
$\Lambda_w$ in the chosen basis.  To check the localization itself,
choose a nonzero integral scalar $d$ with $d z$ in the original global
lattice.  If $z$ has integral coordinates at $w$, choose
$t\in\mathcal O_K\setminus w$ which is divisible, at every other prime
dividing $(d)$, by at least the corresponding valuation of $d$.
Such $t$ exists by the Chinese remainder theorem, with $t\equiv1$
modulo $w$.  At $w$, the coordinates of $tz$ remain integral; elsewhere
$t/d$ is locally integral.  The original \texttt{val0.lat} criterion
now puts $tz$ in the global lattice, so $z\in S_w$.  Conversely a
denominator outside $w$ is a unit in every coordinate in (GLE.17).
This proves the equality without a new lattice choice.

Fix a seen place $w\mid p$ and a selected place $v$ of the original
\texttt{AtTTied} record.  Put $\tau_v=(D.\mathrm{fm}\,u).\mathrm{tauOf}(w,v)$.
Let
\[
 K_{v,c}=\texttt{Lk}\,\mathcal F.\mathrm{atT}\,0\,v\,c
        =\texttt{clsCrd}(c)+P_vH_c,
 \qquad P_v=\mathcal F.\mathrm{atT}.P_v.
 \tag{GLE.18}
\]
Choose a $p$-local basis $B'_{v,c}$ of $K_{v,c}$.  Such bases exist
because localization at the finitely many primes above $p$ is a
semilocal Dedekind ring, hence a PID: every ideal is generated by an
element with its finitely many prescribed valuations.  A finite
torsion-free module over this ring is free; expressing that basis in
the original $E$-space gives the required matrix.  The lattices in
(GLE.18) are full: $H_c$ is a Hermitian dual of a full lattice, and
$\texttt{clsCrd}(c)\subset K_{v,c}\subset H_c$.

Set $C_{v,c}=(B'_{v,c})^{-1}B_c$ at $\tau_v$, and set $C_{v,c}=1$
at the other embeddings.  In the same coordinates the original
neighbor lattice is exactly
\[
 B_{w,v}=(D.\mathrm{fm}\,u).\mathrm{Bw}(0,w,v)
       =\bigoplus_c R_{\lambda_{0,\tau_v}}(C_{v,c})\Lambda_w.
 \tag{GLE.19}
\]
The forward implication is the original \texttt{bw0} clause.
For the reverse implication, each right-hand coordinate is a rational
polynomial by GLE3, since the entries of the basis matrices are in $E$
and the stored embedding is $K$-valued.  Use (GLE.16) to realize those
coordinates in $V$.  Changing either local basis changes the mixed
coordinate matrix by integral invertible matrices at the appropriate
embeddings.  GLE3 preserves $\Lambda_w$ under those changes, so all
the bases quantified in \texttt{bw0} satisfy the required condition.
This proves surjectivity as well as inclusion in (GLE.19).

In particular, if $r$ is the rank of the tensor lattice $\Lambda_w$,
then the literal source dimension is
\[
                  d_S=D.\mathrm{duS}(u)=|\mathcal C|r.
 \tag{GLE.20}
\]
The original \texttt{emb\_basis} identifies scalar extension of $V$
with the complex forms; GLE3 and (GLE.16) calculate its dimension.
No identification of $d_S$ with the smaller ramified dimension is made.
% END GLE4

\subsection{Integral maps and the second branch on those lattices}

% BEGIN GLE5
\paragraph{GLE5: the two original Hecke operators are integral.}
\label{claim:GLE5}
Supply the same original $\mathcal F$, its finite level character with
values in $K$, and \texttt{admS}.  For each seen $w\mid p$ and every
selected $v$, the original operators satisfy
\[
              U_v(B_{w,v})\subseteq S_w,
              D_v(S_w)\subseteq B_{w,v}.           \tag{GLE.21}
\]
They are $K$-rational operators on $V$.  This is the complete literal
\texttt{HeckeIntegralLevelZero}; it in particular proves the
source's version with the additional \texttt{HidaSetting} hypothesis.

We verify the geometric and coefficient parts separately on the same
move.  Work at $P_v$ and use the original residue map.  We first check
the precise uniformizer assertion, following the elementary norm
argument of FC1.  There are $f$ distinct selected primes, each of
residue cardinality $p$, and every one divides $\pi_u$.  The original
host has degree $[E:\mathbb Q]=2f$ and
$\pi_u\bar\pi_u=p$.  Thus the absolute ideal norm of $(\pi_u)$ is
$p^f$: its square is the norm $p^{2f}$ of $(p)$, and conjugation
preserves ideal norm.  The $f$ displayed primes already account for
this whole norm.  Each therefore occurs in $(\pi_u)$ with exponent
one.  The original nonzero residue of $\bar\pi_u$ makes it a unit at
each of them.  Consequently $p$ has valuation one at $P_v$, and
$P_v\mathcal O_{E,P_v}=p\mathcal O_{E,P_v}$.
Write $H$ for the localized $H_c$, and let $z$ lift an up index with
$f(z)=1$.  The exact up naturality fields, already reconstructed in
NB2 and NB7, say that its rational unitary transport $\gamma$ identifies
\[
 \gamma\,\texttt{clsCrd}(c^+)=H,\qquad
 \gamma H_{c^+}=H+\mathcal O_{E,P_v}p^{-1}z.
 \tag{GLE.22}
\]
Consequently
\[
 \gamma (K_{v,c^+})_{P_v}
   =H+p\bigl(H+\mathcal O_{E,P_v}p^{-1}z\bigr)=H.
 \tag{GLE.23}
\]
For a down index $g$ with $g(t_e)=1$, NB3 and NB8 identify the original
transport by
\[
             \gamma(H_{c^-})_{P_v}
                     =(\texttt{clsCrd}(c))_{P_v}=(K_{v,c})_{P_v}.
 \tag{GLE.24}
\]
For the last equality, $pH_c\subseteq\texttt{clsCrd}(c)$ follows from
the original base-lattice dual bound.  At each selected place $v'\ne v$
the transported dual lattices are equal, because the adelic move is
the identity there and the residual level factor preserves the base
lattice and its Hermitian dual.  To descend that completed equality
to rational localization, clear the finitely many denominators at the
moved place with an element which is a unit at $P_{v'}$.  Equivalently,
the two rational lattices have identical localizations away from the
support of their finite quotient.  This proves the exact
\texttt{AwayLat0} assertion used by the source.

Let $B_{\rm in}$ denote the appropriate incoming base or neighbor
basis, and $B_{\rm out}$ the outgoing one.  Equations (GLE.23)--(GLE.24)
and their away-place versions say precisely that
\[
       B_{\rm in}^{-1}\gamma^{-1}B_{\rm out}
                \in\mathrm{GL}_k(\mathcal O_{E,P_{v'}})
 \tag{GLE.25}
\]
at the prime indexed by each $v'$.  Its inverse is integral too.
Here the place dictionary on the CM field deserves a check:
\texttt{tauOf} states a residue equality on the totally real field,
whereas these matrices have entries in $E$.  Let $\nu$ be the
$w$-valuation after the stored embedding $E\to K$ at
$\tau=(D.\mathrm{fm}\,u).\mathrm{tauOf}(w,v')$.  It is nonnegative on
$\mathcal O_E$.  The original \texttt{seenP} gives $\nu(\pi_u)>0$.
The element $\pi_u+\bar\pi_u$ lies in the totally real field and has
nonzero residue at $P_{v'}$.  Thus \texttt{tauOf} gives it valuation
zero, and hence $\nu(\bar\pi_u)=0$.  If $y\in P_{v'}$, then
$y\bar\pi_u+\bar y\pi_u$ is totally real and has zero residue.
Its valuation is positive, as is that of $\bar y\pi_u$; subtraction
shows $\nu(y\bar\pi_u)>0$, so $\nu(y)>0$.
If $b\notin P_{v'}$, choose $a\in\mathcal O_E$ with
$ab-1\in P_{v'}$.  Were $\nu(b)>0$, both $ab$ and $ab-1$ would have
positive valuation, a contradiction.  Thus $\nu(b)=0$.
It follows that $\mathcal O_{E,P_{v'}}$ maps to $\mathcal O_w$,
with every denominator outside $P_{v'}$ mapping to a unit.  Applying
this to the matrix and its inverse in (GLE.25) gives an integral
invertible coefficient matrix at every $\tau$.

The original function transport is
$\chi(\kappa)^{-1}$ times substitution by $\gamma^{-1}$ for the
source's factorization $\gamma y_{c'}=y_c\eta\kappa$, including
the determinant powers in (GLE.14).  Its character factor is a unit
at every coefficient place: if $\chi(\kappa)=\iota_K(r)$ and
$\chi(\kappa)^q=1$ with $q>0$, injectivity gives $r^q=1$;
both $r$ and $r^{-1}$ are algebraic integers.  GLE3 applied to
(GLE.25) therefore sends each integral incoming polynomial to an
integral outgoing polynomial.  This is valid for every choice of
local bases, so GLE4 identifies the resulting conditions with the
actual lattices in (GLE.21).

Finally the operators are the original finite sums
\[
 (U_v\phi)(c)=\sum_{z:f_v(z)=1}\operatorname{upT}_{v,c,z}
                         \phi(c^+_z),\qquad
 (D_v\phi)(c)=\sum_{g:g(t_e)=1}\operatorname{dnT}_{v,c,g}
                         \phi(c^-_g).             \tag{GLE.26}
\]
There is no averaging denominator.  Naturality identifies these sums
with the finite adelic translates, and changing a representative
permutes their residue indices.  They consequently remain in the
same complex forms space.  Every summand has a $K$-rational
polynomial value: $\gamma$ is over $E$, its stored embeddings take
values in $K$, and the character factor is in $K$.  The reverse
direction of the original \texttt{val0.rat} condition therefore places
the sum in the original $V$.  Summing the integral coordinate
conditions proves both inclusions in (GLE.21), with exactly the
source's \texttt{Into} preimages under its injective embedding.
% END GLE5

% BEGIN GLE6
\paragraph{GLE6: the exact second-branch inclusion and length.}
\label{claim:GLE6}
For the same original forms and trivial stabilizers, fix seen $w\mid p$
and $v$, and put
$b_v=(D.\mathrm{fm}\,u).bTop(\tau_v)=-\lambda_{0,\tau_v}(k-1)$.
If $b_v\ge0$, then
\[
 \begin{split}
 p^{b_v}S_w&\subseteq B_{w,v},\\
 \ell_{\mathcal O_w}(B_{w,v}/p^{b_v}S_w)
   &=\frac{e_wd_S}{k}
           \bigl(|\lambda_{0,\tau_v}|+kb_v\bigr).
 \end{split}                                      \tag{GLE.27}
\]
In particular the displayed rational number is a nonnegative integer.
This is the full original \texttt{SecondBranchInclusionLength}.

The original residue functional on $H_c$ is nonzero, has scalar
residue-field target of cardinality $p$, and has kernel $K_{v,c}$.
This is exactly \texttt{lv0.fv\_ker}; surjectivity follows by scaling
one nonzero value.  Thus, over $\mathcal O_{E,P_v}$, this is a
colength-one sublattice of $H_c$.  Choosing a lift $z_0$ with
$f(z_0)=1$ and lifting a basis of its kernel gives bases with relative
matrix $\operatorname{diag}(p,1,\ldots,1)$.  The prime $p$ has valuation
one here, as in GLE5.  Any chosen global $p$-local bases have the same
Smith form after localization.  Under the original embedding at
$\tau_v$, their relative matrix therefore has the form
\[
 B_c^{-1}B'_{v,c}
       =A_c\operatorname{diag}(p,1,\ldots,1)C_c,
       \qquad A_c,C_c\in\mathrm{GL}_k(\mathcal O_w).
 \tag{GLE.28}
\]
This description is made only in the coefficient DVR for computing
the existing lattices; it does not replace their global bases.
The other factors in the tensor representation are unchanged.

By GLE3, the unit matrices in (GLE.28) preserve $\Lambda_w$, and
also preserve its scalar multiple $p^{b_v}\Lambda_w$.
It suffices to calculate the quotient with the inverse diagonal
matrix in (GLE.19).  In a homogeneous row-degree basis, its entries
are $p^{-\beta_h}$, where
$\beta_h=d_1(h)+\lambda_{0,\tau_v}(k-1)$.
For every such basis element,
$b_v+\beta_h=d_1(h)\ge0$.  Thus the inclusion holds and its elementary
divisor has length $e_w(b_v+\beta_h)$.  Tensor the basis with the
bases of the other factors and use (GLE.15).  On one class the length
is
\[
                  e_wr\bigl(b_v+|\lambda_{0,\tau_v}|/k\bigr).
\]
Sum over the same finite class set and use the exact identity
$d_S=|\mathcal C|r$ from (GLE.20).  This proves (GLE.27), including
finiteness and the exact normalization by $e_w$ and $k$.
The determinant shift is essential: omitting it would give the wrong
index when the lowest weight is negative.
% END GLE6

\subsection{The finite local calculation on the original ramified forms}

% BEGIN GLE7
\paragraph{GLE7: diamond matchings and the exact local move dictionary.}
\label{claim:GLE7}
Fix a selected place $v$, write $\mathbb F=\mathbb F_p$ for its
original residue field, and localize the original lattices at $P_v$.
Write $H$ for the dual lattice, $L$ for the primal lattice, and
$f:H/pH\to\mathbb F$ for the original nonzero functional.  The
source's \texttt{fv\_ker} identifies $L/pH$ with $\ker f$, so
$pH\subset L\subset H$.  Put $V_0=L/pL$, of dimension $k$, and
choose $z_0\in H$ with $f(z_0)=1$.  Its original distinguished vector is
\[
                         t_e=pz_0\pmod {pL}\ne0.                 \tag{GLE.30}
\]
This is independent of $z_0$: two such lifts differ by an element of
$L$.  Nonvanishing follows because $pz_0\in pL$ would put $z_0\in L$.
The prescribed residue projections identify this $V_0$ and $t_e$
with \texttt{lk} and \texttt{te}, without a choice of a new residue
space.

The original local moves have the following precise description.
An up index is $z\in H/pH$ with $f(z)=1$.  After its original rational
transport the new dual lattice is $H+p^{-1}\mathcal O z$, the new
primal lattice is $H$, and the new functional is
\[
                         f^+(y)=f(py).
\]
An arbitrary down index is a nonzero functional $g\in V_0^*$.
After its original transport the new dual lattice is $L$, its new
primal lattice is the inverse image of $\ker g$ in $L$, and the new
functional is $g$.  The ordinary down operator restricts these
indices by $g(t_e)=1$; the auxiliary operator $E_v$ uses all
$g\ne0$.  These are the local equalities of NB2--NB3 and NB7--NB8.
At the conjugate prime the data are fixed by Hermitian duality;
at every other prime they are unchanged.

The original diamond indexed by $a\in\mathbb F^\times$ replaces
$f$ by $af$ and leaves $H,L$ unchanged.  It replaces $t_e$ by
$a^{-1}t_e$.  In an up sum the bijection $z\mapsto a^{-1}z$
preserves the enlarged lattice and multiplies the resulting
functional by $a$.  In a down sum the bijection $g\mapsto ag$
preserves the condition $g(t_e)=1$ after this change of $t_e$, and
again multiplies the resulting functional by $a$.
These give, on the original forms,
\[
 \rho_{v,a}U_v=U_v\rho_{v,a},\qquad
 \rho_{v,a}D_v=D_v\rho_{v,a},\qquad E_v\rho_{v,a}=E_v.          \tag{GLE.31}
\]
For the last equality, all nonzero functionals occur, and multiplication
by $a$ simply permutes them.  At two distinct selected places the
local moves have disjoint support.  Their finite-adelic matrices
commute componentwise, which proves the corresponding commutations
of $U_v,D_v$ with every $\rho_{v',a}$ for $v'\ne v$.

We check why these identifications are equalities of transported
values, rather than only equalities of lattice shapes.  Work at the
original preferred adelic representative and use moves supported at
the place in question.  If their final dual lattice, primal lattice
and functional agree, the ratio of the two move matrices stabilizes
these data.  The original \texttt{UlocAlt} and \texttt{U\_eq}
identify this stabilizer with the original level subgroup.  Its
components away from the moved place are the identity.  The selected
place is disjoint from the two auxiliary bad primes: RSL8 proves this
from $\pi_u\equiv1$ modulo $(g_0g_1)^2$ and
$\pi_u\bar\pi_u=p$.  The original \texttt{chi\_triv} therefore
makes the level character of the ratio equal to one.  Adelic
equivariance gives equal values, with no residual scalar.
Factoring the moves through the prescribed class representatives
introduces the same natural transforms as in (GLE.26), so it does
not alter the equality.  This also proves preservation of the
original forms space by $U_v,D_v,E_v$ and the diamonds: changing a
level representative permutes the displayed finite indices and
uses this same equivariance.  No transitivity or approximation
assertion for a different class space is being invoked.
% END GLE7

% BEGIN GLE8
\paragraph{GLE8: the actual $DU$ formula by finite fibers.}
\label{claim:GLE8}
Continue with exactly the preceding local data.  First make a down
move indexed by $g\ne0$, then an up move from its new dual lattice
$L$.  The original down transport identifies the latter residue
space with $V_0=L/pL$.  Thus its up indices are in bijection with
\[
                         \{y\in V_0:g(y)=1\}.                  \tag{GLE.32}
\]
Both directions use the original residue projections: a lift of $y$
in $L$ pulls back to a lift in the down dual lattice, and two lifts
give the same index exactly when their difference lies in $pL$.
The projection is surjective and has this kernel by \texttt{prk}
and \texttt{prk\_ker}.  This verifies the whole bijection, including
its uniqueness.

For a nonzero $y\in V_0$ choose a lift $\widetilde y\in L$.
After the two transports the final dual and primal lattices and
functional are
\[
 H_y=L+p^{-1}\mathcal O\widetilde y,\qquad L_y=L,
 \qquad h_y(l+c p^{-1}\widetilde y)=\bar c.                    \tag{GLE.33}
\]
Here $\bar c$ denotes reduction to $\mathbb F$.
Indeed the up functional is $g(p\,\cdot)$, and $g(y)=1$;
its kernel in $H_y/pH_y$ is the image of $L$.
The expression is well defined: if $c p^{-1}\widetilde y\in L$,
then $\bar c\,y=0$, hence $\bar c=0$.
Changing the lift by $pL$ does not change any of these data.
In particular the endpoint depends only on $y$, not on $g$.
The conjugate local component is determined by its Hermitian dual,
and all other components are unchanged.  The character-one
comparison in GLE7 therefore proves equality of the actual
transported form values.  Denote that common value, at a fixed class
and frame, by $W(y)$.

For $y=t_e$, (GLE.33) gives the original $H,L,f$, so $W(t_e)=\phi$.
For $y=a t_e$ with $a\ne0$ it gives $H,L,a^{-1}f$.  Thus
$W(a t_e)=\rho_{v,a^{-1}}\phi$ at that class and frame.  In
particular
\[
             \sum_{a\in\mathbb F^\times} W(a t_e)
                       =\operatorname{Nm}_v\phi,\qquad
             \operatorname{Nm}_v=\sum_a\rho_{v,a}.             \tag{GLE.34}
\]
The inversion in this formula records the literal source diamond
convention.

It remains to count fibers, a calculation in the $k$-dimensional
space $V_0$.  If $y\ne0$, choosing a basis beginning with $y$
shows that exactly $p^{k-1}$ functionals satisfy $g(y)=1$;
none of them is zero.  For $y=0$ there are none.
Requiring in addition $g(t_e)=1$ gives
\[
 \#\{g:g(t_e)=g(y)=1\}=
 \begin{cases}
 p^{k-1},&y=t_e,\\
 0,&y\in\mathbb F t_e, y\ne t_e,\\
 p^{k-2},&y\notin\mathbb F t_e.
 \end{cases}                                                  \tag{GLE.35}
\]
The middle case follows by linearity.  In the last case extend
$t_e,y$ to a basis; two coordinates of $g$ are fixed and the other
$k-2$ are arbitrary.  The original ranks have $k\ge5$, so the
exponents are ordinary nonnegative integers.

Expanding the exact finite sums, reindexing by (GLE.32), and using
(GLE.35) now gives
\[
 \begin{split}
 D_vU_v\phi&=p^{k-1}\phi+
                   p^{k-2}\sum_{y\notin\mathbb F t_e}W(y),\\
 E_vU_v\phi&=p^{k-1}\sum_{y\ne0}W(y).
 \end{split}
\]
Subtract the line contribution (GLE.34).  With the original
normalization $\operatorname{trp}_v=p^{-(k-1)}E_vU_v$, the result is
\[
 D_vU_v\phi=p^{k-1}\phi-p^{k-2}\operatorname{Nm}_v\phi
                              +p^{k-2}\operatorname{trp}_v\phi. \tag{GLE.36}
\]
Every step holds at each original class and frame, so this is an
identity in the original forms space.  It is the complete literal
\texttt{DUFormulaLevelZero}; its proof has supplied the move,
fiber, endpoint and scalar identifications needed for that formula.
% END GLE8

% BEGIN GLE9
\paragraph{GLE9: the inverse composite on the original ramified space.}
\label{claim:GLE9}
Let
\[
 M^{\rm ram}=\mathcal F.\mathrm{Fi}(0)\cap
                  \bigcap_{v'}\ker\operatorname{Nm}_{v'}
                    =\texttt{MramI D u 0}.
\]
The commutations (GLE.31), including those at distinct places,
show that $U_v$ and $D_v$ preserve every displayed kernel and the
forms space.  Moreover (GLE.31) gives
\[
 \operatorname{trp}_v\rho_{v,a}
       =p^{-(k-1)}E_vU_v\rho_{v,a}
       =p^{-(k-1)}E_v\rho_{v,a}U_v
       =\operatorname{trp}_v.
\]
For $\phi\in M^{\rm ram}$ this implies
\[
 0=\operatorname{trp}_v(\operatorname{Nm}_v\phi)
                 =(p-1)\operatorname{trp}_v\phi.
\]
Since $p$ is a prime, $p-1\ne0$ in $\mathbb C$, and hence
$\operatorname{trp}_v\phi=0$.  Equation (GLE.36) reduces to
$D_vU_v\phi=p^{k-1}\phi$.

The original forms space is finite dimensional by its original
\texttt{emb\_basis}, so its subspace $M^{\rm ram}$ is too.
The last identity makes $U_v$ injective there.  A linear injective
endomorphism of a finite-dimensional space is surjective, and
therefore $D_v=p^{k-1}U_v^{-1}$ on this same space.  Consequently
\[
 U_v(M^{\rm ram})\subseteq M^{\rm ram},\qquad
 D_v(M^{\rm ram})\subseteq M^{\rm ram},\qquad
 U_vD_v\phi=p^{k-1}\phi\quad(\phi\in M^{\rm ram}).             \tag{GLE.37}
\]
This is exactly \texttt{HeckeRamifiedLevelZero}, with $k=N(0)$.
No change of the ramified subspace, of the finite character, or of
the preferred representative system has occurred.  The proof uses
the original \texttt{Forms0} ties; neither a new abstract Hecke
module nor an additional matching predicate is required.
% END GLE9

\subsection{The original congruence lattice and its length}

% BEGIN GLE10
\paragraph{GLE10: a DVR estimate for one normalized original operator.}
\label{claim:GLE10}
Let $\mathcal O=\mathcal O_w$, let $e=e_w$ be the valuation of $p$,
and write $S=S_w$, $B=B_{w,v}$, $W=\mathcal F.\mathrm{ramSub}$.
The rationality in GLE5 descends the two original operators uniquely
to $K$-linear operators on $V$.  Their descent is unique because
the original embedding is injective.  GLE9 and the original
\texttt{ramSub} equality make both descend to $W$, with
$UD=p^{k-1}I$ there.  Put
\[
 R=S\cap W,\qquad B_R=B\cap W,\qquad b=b_v\ge0,
 \qquad G=p^bU.
\]
Thus $R$ is the original \texttt{Sram}, not a substitute lattice.
The equalities $p^bR=(p^bS)\cap W$ and
$U(p^bR)=p^bU(R)$ follow because $W$ is a $K$-subspace.
GLE5--GLE6 and GLE9 give
\[
 p^bR\subseteq B_R,\qquad U(B_R)\subseteq R,\qquad
 D(R)\subseteq B_R,\qquad
 p^{k-1}R\subseteq U(B_R).
\]
In particular $G(R)\subseteq R$.

All these are finitely generated torsion-free modules over a DVR,
hence free.  Multiplication by $p$ has cokernel length $e$ on each
free generator.  The last inclusion gives
\[
 \ell(R/U(B_R))\le e(k-1)\operatorname{rank}_{\mathcal O}R
                                       \le e(k-1)d_S.
\]
The natural injection
$B_R/p^bR\hookrightarrow B/p^bS$ shows
$\ell(B_R/p^bR)\le\ell(B/p^bS)$.
Applying $U$ gives a surjection
$B_R/p^bR\twoheadrightarrow U(B_R)/U(p^bR)$, so the same upper
bound holds for the latter quotient.  The exact sequence of the
three nested modules
$U(p^bR)\subseteq U(B_R)\subseteq R$ therefore yields
\[
 \ell(R/G(R))\le
      \ell(B/p^bS)+e(k-1)d_S
   =e d_S\left(\frac{|\lambda_{0,\tau_v}|+kb_v}{k}+k-1\right).
                                                               \tag{GLE.38}
\]
The term on the left is finite; each step has just exhibited a
finite-length quotient that bounds it.  The rank inequality uses
only $\operatorname{rank}R\le\dim_K V=d_S$, established on the
original coordinates in GLE4.  This independently proves the
source's local \texttt{LemmaC\_core} estimate with the exact
original dimension and normalization.
% END GLE10

% BEGIN GLE11
\paragraph{GLE11: the ordered product, its power, and the complete length bound.}
\label{claim:GLE11}
Assume \texttt{HidaSetting D u} and fix $w\mid p$.
Let $\mathcal V_w$ be the original \texttt{VVT}: it consists of the
selected places whose $\tau_v$ is top if $w$ is seen, and is empty
otherwise.  On this set $b_v\ge0$ by the last clause of
\texttt{HidaSetting}.  For seen $w$ let $G_v=p^{b_v}U_v$ on the
same $R=\texttt{Sram D u w}$.  GLE10 applies to every $v$.

We spell out the finite-index argument to avoid any commutativity
assumption about the chosen order.  If $A(R),B(R)\subseteq R$, then
$A(R)/AB(R)$ is a quotient of $R/B(R)$ under $A$.  The exact
sequence with submodule $A(R)/AB(R)$ gives
\[
 \ell(R/AB(R))\le\ell(R/A(R))+\ell(R/B(R)).                    \tag{GLE.39}
\]
Induct over the original \texttt{toList.foldr} order and then over
the $m$ powers.  The resulting operator $G_*$ satisfies
\[
 G_*^m(R)\subseteq R,\qquad
 \ell(R/G_*^m(R))\le
 m e_wd_S\sum_{v\in\mathcal V_w}
       \left(\frac{|\lambda_{0,\tau_v}|+kb_v}{k}+k-1\right).
                                                               \tag{GLE.40}
\]
Scalars commute with all linear maps, so, in that very same order,
\[
 G_*^m=p^{m\sum_{v\in\mathcal V_w}b_v}
       \left(\prod_{v\in\mathcal V_w}^{\rm original\ order}U_v\right)^m.
\]
The exponent and product are literally \texttt{expT} and
\texttt{Uprod}.  The source's definition of \texttt{imgL}, together
with the unique rational descent in GLE10, therefore identifies
\[
              G_*^m(R)=\texttt{Nlam D u w}.                    \tag{GLE.41}
\]
This is equality of submodules of the original rational $V$.

For completeness, the weight simplification in the stated source
bound is also elementary.  At a top $\tau$, the first entry is
at most $a_{\max}$, and the absolute values of every later entry
are at most the source's $\operatorname{Wt}(\tau)$.  The last
entry has absolute value $b_v\le\operatorname{Wt}(\tau)$.
Thus
\[
 \frac{|\lambda_{0,\tau}|+kb_v}{k}
 \le \frac{a_{\max}}k+
       \frac{k-1}{k}\operatorname{Wt}(\tau)+b_v
 \le \frac{a_{\max}}k+2\operatorname{Wt}(\tau).
\]
Here $|\lambda|$ continues to denote the signed sum of its entries,
not their absolute-value sum.  Using $k-1\le k$ and $b_v\ge0$ in
(GLE.40) proves the exact, slightly looser source bound
\[
 \begin{split}
 \texttt{Nlam D u w}&\subseteq\texttt{Sram D u w},\\
 \ell(\texttt{Sram}/\texttt{Nlam})
 &\le e_w m d_S
   \sum_{v\in\mathcal V_w}
     \left(\frac{a_{\max}}k+2\operatorname{Wt}(\tau_v)+k+b_v\right).
 \end{split}                                                  \tag{GLE.42}
\]
The length is a natural number, as proved in (GLE.40).
If $w$ is unseen, the original list is empty, its product is the
identity, and its exponent is zero.  Hence \texttt{Nlam} equals
\texttt{Sram} and both sides of the required length inequality are
zero.  This handles that case without presupposing a seen place.
Equations (GLE.41)--(GLE.42) are the complete literal
\texttt{CongruenceLatticeLengthBound D u}.
% END GLE11

% BEGIN GLE12
\paragraph{GLE12: uniform installation and the remaining common producer.}
\label{claim:GLE12}
The logical roles of the preceding statements are distinct.
GLE1--GLE2 prove
\[
             \texttt{TuSize}(\mathcal D)
                 \Longrightarrow\texttt{SettingEventually}(\mathcal D)
                                                               \tag{GLE.43}
\]
for the actual family with its original parameter witnesses.
GLE5, GLE9 and GLE6 respectively prove the literal pointwise
\texttt{HeckeIntegralLevelZero}, \texttt{HeckeRamifiedLevelZero},
and \texttt{SecondBranchInclusionLength} from the same original
\texttt{Forms0}, its finite $K$-valued level character and
\texttt{admS}.  The second-branch statement retains its explicit
$b_v\ge0$ hypothesis.  GLE11 consumes these three statements and
\texttt{HidaSetting}, on the same $D,u$, to produce the complete
congruence lattice length bound.

In the selected application the original
\texttt{RealPlaceSecondLayerConditions} supplies, eventually, the
same \texttt{FormsTiedSym}, its \texttt{ChiFin} data, and the
required real-place ties.  RSL6 fixes the common original selector;
RSL7 supplies \texttt{admS} on those forms.  Take the maximum of
that supplier threshold and the threshold in GLE2.  It is chosen
before $i$, before the word $u$, and before the places $w,v$.
The pointwise proofs above are uniform in all of these variables.
Thus the original guarded family predicate is satisfied for each
of the three operator conclusions:
\[
 \exists A\ \forall x\ge A\ \forall i\ \forall u,
 \quad \mathcal D.x_1<x\ \Longrightarrow
 \bigl(\texttt{HidaSetting}(D_{x,i},u)\Longrightarrow
          P(D_{x,i},u)\bigr).                                  \tag{GLE.44}
\]
This is the literal \texttt{FamU} guard and quantifier order.
One may then discharge its guard using (GLE.43), on the same late
family.  Maxima of finitely many thresholds do not change the
family, its class constants, any parameter witness, or the
selector.  Accordingly this also transports to the source's
\texttt{Closed} predicates whenever those upstream predicates are
supplied on a class.  It does not choose a fresh family separately
for the three conclusions.

The five implication tasks
\texttt{TASK\_GL6\_SIZE\_SETTING},
\texttt{TASK\_GL6\_INTEGRAL},
\texttt{TASK\_GL6\_RAMIFIED},
\texttt{TASK\_GL6\_SECOND\_BRANCH}, and
\texttt{TASK\_GL6\_CONGRUENCE\_LENGTH} are thereby closed in this
stated scope.  The first now composes with the original construction
class's \texttt{TuSize} producer; the other four retain the
original same-selector forms and finite-character input as explicit
upstream columns.  The compatible common selector remains a
producer obligation, and this chapter does not assert the GL4
$hZ_a,hZ_b,t_{CV}$ inputs or the entire eleven-field
\texttt{FamGLmin} output.  Once their separate inputs are supplied,
the present conclusions are available on exactly the original
families required for that final assembly.  RSL11--RSL12 may also
reuse the inclusion in (GLE.42) to obtain their already proved
fullness conclusion without a separate congruence-lattice inclusion
producer.
% END GLE12


